\Article{3}
If $m$ denotes a complex integer having no common divisor with $\mu$, and every number of a quarter system $\kappa$ is multiplied by $m$, the products, or any numbers congruent to them, themselves form a quarter system. Likewise, multiplying the numbers of the systems $\kappa',\kappa'',\kappa'''$ gives three further quarter systems associated with the first. The proof is easy to give. Denote these four new systems by $m\kappa,m\kappa',m\kappa'',m\kappa'''$, whether the products themselves or only numbers congruent to them are chosen. In the latter case, the choice can always be made from one of the systems $\kappa,\kappa',\kappa'',\kappa'''$. Thus the system $\chi$ has been partitioned into quarter systems in two different ways, at least generally speaking. Suppose that $m\kappa$ has in common

\[
\begin{array}{rl}
\text{with }\kappa\dotsfive&\lambda\text{ numbers}\\
\kappa'\dotsfive&\lambda'\text{ numbers}\\
\kappa''\dotsfive&\lambda''\text{ numbers}\\
\kappa'''\dotsfive&\lambda'''\text{ numbers}
\end{array}
\]
Then $\kappa'$ with $m\kappa'$, $\kappa''$ with $m\kappa''$, and $\kappa'''$ with $m\kappa'''$ will also have $\lambda$ members in common;
\SourceLeaf{380}
$\kappa''$ with $m\kappa'$, $\kappa'''$ with $m\kappa''$, $\kappa$ with $m\kappa'''$, $\lambda'$ members, and so on. Let $\varepsilon$ be the least residue of $\lambda'+2\lambda''+3\lambda'''$ modulo $4$, or let $\varepsilon$ be one of the four numbers $0,1,2,3$ according as $\lambda'+2\lambda''+3\lambda'''$ has the form $4n,4n+1,4n+2,4n+3$. I now claim that $\varepsilon$ is independent of the arrangement of the quarter system $\kappa$.

To make the proof easier, I use the following notation. Let $\Pi\psi$ denote the number $0,1,2,3$ according as the number $\psi$, not divisible by $\mu$, occurs (itself or through representation by congruence) in the group $\kappa,\kappa',\kappa'',\kappa'''$. The following consequences are therefore immediate:

I.\quad $\Pi(i\psi)\equiv1+\Pi\psi\ (\text{mod. }4)$.

II.\quad The number $i^{-\Pi\psi}\pdot\psi$ occurs, itself or through representation by congruence, in the group $\kappa$.

III.\quad $\Sigma\Pi m\varphi\equiv\varepsilon,\ (\text{mod. }4)$, when the sum extends over all members $\varphi$ occurring in $\kappa$.

Now let $k$ be another quarter system, consisting of $f,f',f''$, and so on, while $\kappa$ consists of $\varphi,\varphi',\varphi''$, and so on. I assume, as is allowed because the \emph{order} of the members of $\kappa$ is arbitrary, that $f$ is identical with or associated with $\varphi$, $f'$ with $\varphi'$, $f''$ with $\varphi''$, and so on. Let the quarter groups associated with $k$ be $k'(\equiv ik)$, $k''(\equiv-k)$, $k'''(\equiv-ik)$. Furthermore, let the characteristic $\mathrm{P}$ have the same meaning for the groups $k,k',k'',k'''$ as $\Pi$ has for $\kappa,\kappa',\kappa'',\kappa'''$, so that $\mathrm{P}\psi=0,1,2,3$ according as $\psi$ belongs to $k,k',k'',k'''$.

Thus, if the partition of $\chi$ into the quarter systems $k,k',k'',k'''$ is used instead of $\kappa,\kappa',\kappa'',\kappa'''$, the least residue of $\mathrm{P}mf+\mathrm{P}mf'+\mathrm{P}mf''+\mathrm{P}mf'''$, and so on, and so on, or of $\Sigma\mathrm{P}mf$, modulo $4$, takes the place of $\varepsilon$; and it remains to prove that $\Sigma\mathrm{P}mf-\Sigma\Pi m\varphi$ is divisible by $4$.

We write this quantity as
\[
\begin{array}{l}
\phantom{+}\mathrm{P}mf+\mathrm{P}mf'+\mathrm{P}mf''+\mathrm{P}mf'''+\text{and so on}\\
-\mathrm{P}m\varphi-\mathrm{P}m\varphi'-\mathrm{P}m\varphi''-\mathrm{P}m\varphi'''-\\
+\mathrm{P}m\varphi+\mathrm{P}m\varphi'+\mathrm{P}m\varphi''+\mathrm{P}m\varphi'''+\\
-\Pi m\varphi-\Pi m\varphi'-\Pi m\varphi''-\Pi m\varphi'''-
\end{array}
\]

Since, by assumption, $f$ and $\varphi$ are congruent or associated, the same holds for $mf$ and $m\psi$, and

\SourceLeaf{381}
\[
\left.\begin{aligned}
f&\equiv i^{-\mathrm{P}\varphi}\varphi\\
i^{-\mathrm{P}mf}mf&\equiv i^{-\mathrm{P}m\varphi}m\varphi
\end{aligned}\right\}\ \text{mod. }\mu
\]
From this it follows easily that $\mathrm{P}mf-\mathrm{P}m\varphi\equiv-\mathrm{P}\varphi\ (\text{mod. }4)$ and that the sum of the top two rows $\equiv-\Sigma\mathrm{P}\varphi$. Furthermore, since $\mathrm{P}m\varphi-\Pi m\varphi\equiv\mathrm{P}(i^{-\Pi m\varphi}m\varphi)$, $m\varphi\pdot i^{-\Pi m\varphi}$ belongs to $\kappa$, and the collection of \emph{all} $m\varphi\pdot i^{-\Pi m\varphi}$, regardless of order, agrees with all $\varphi$, the sum of \emph{all} $\mathrm{P}(m\varphi\pdot i^{-\Pi m\varphi})$ will equal the sum of all $\mathrm{P}\varphi$. Thus the sum of the third and fourth rows $\equiv\Sigma\mathrm{P}\varphi\ (\text{mod. }4)$, and hence the sum of all four rows $\equiv0\ (\text{mod. }4)$. Which was to be proved.

Since $\varepsilon$ thus depends only on $m$ and $\mu$, independently of the choice of quarter systems, we shall call $\varepsilon$ the character of the number $m$ in relation to the modulus $\mu$ and denote it by $\text{Ch. }m\ (\text{mod. }\mu)$. It is easy to see that this only generalizes, or includes, the definition given (Art. $\ldots$) for the case in which $\mu$ is prime.
